% Options for packages loaded elsewhere
\PassOptionsToPackage{unicode}{hyperref}
\PassOptionsToPackage{hyphens}{url}
\PassOptionsToPackage{dvipsnames,svgnames,x11names}{xcolor}
\documentclass[
  11pt,
]{article}
\usepackage{xcolor}
\usepackage[margin=25mm]{geometry}
\usepackage{amsmath,amssymb}
\setcounter{secnumdepth}{-\maxdimen} % remove section numbering
\usepackage{iftex}
\ifPDFTeX
  \usepackage[T1]{fontenc}
  \usepackage[utf8]{inputenc}
  \usepackage{textcomp} % provide euro and other symbols
\else % if luatex or xetex
  \usepackage{unicode-math} % this also loads fontspec
  \defaultfontfeatures{Scale=MatchLowercase}
  \defaultfontfeatures[\rmfamily]{Ligatures=TeX,Scale=1}
\fi
\usepackage{lmodern}
\ifPDFTeX\else
  % xetex/luatex font selection
\fi
% Use upquote if available, for straight quotes in verbatim environments
\IfFileExists{upquote.sty}{\usepackage{upquote}}{}
\IfFileExists{microtype.sty}{% use microtype if available
  \usepackage[]{microtype}
  \UseMicrotypeSet[protrusion]{basicmath} % disable protrusion for tt fonts
}{}
\makeatletter
\@ifundefined{KOMAClassName}{% if non-KOMA class
  \IfFileExists{parskip.sty}{%
    \usepackage{parskip}
  }{% else
    \setlength{\parindent}{0pt}
    \setlength{\parskip}{6pt plus 2pt minus 1pt}}
}{% if KOMA class
  \KOMAoptions{parskip=half}}
\makeatother
\usepackage{listings}
\newcommand{\passthrough}[1]{#1}
\lstset{defaultdialect=[5.3]Lua}
\lstset{defaultdialect=[x86masm]Assembler}
\usepackage{longtable,booktabs,array}
\usepackage{calc} % for calculating minipage widths
% Correct order of tables after \paragraph or \subparagraph
\usepackage{etoolbox}
\makeatletter
\patchcmd\longtable{\par}{\if@noskipsec\mbox{}\fi\par}{}{}
\makeatother
% Allow footnotes in longtable head/foot
\IfFileExists{footnotehyper.sty}{\usepackage{footnotehyper}}{\usepackage{footnote}}
\makesavenoteenv{longtable}
\setlength{\emergencystretch}{3em} % prevent overfull lines
\providecommand{\tightlist}{%
  \setlength{\itemsep}{0pt}\setlength{\parskip}{0pt}}
\emergencystretch=3em
\AtBeginDocument{\Urlmuskip=0mu plus 2mu}
\lstset{breaklines=true,basicstyle=\small\ttfamily,columns=fullflexible,keepspaces=true}
\setlength{\tabcolsep}{4pt}
\usepackage{iftex}
\ifXeTeX
  \usepackage{newunicodechar}
  \makeatletter
  \newcommand{\DeclareUnicodeCharacter}[2]{%
    \begingroup
    \uccode`\~="#1\relax
    \uppercase{\def\temp{~}}%
    \expandafter\endgroup
    \expandafter\newunicodechar\expandafter{\temp}{#2}%
  }
  \makeatother
\fi
\ifdefined\DeclareUnicodeCharacter
\DeclareUnicodeCharacter{2192}{\ensuremath{\to}}
\DeclareUnicodeCharacter{2190}{\ensuremath{\leftarrow}}
\DeclareUnicodeCharacter{2194}{\ensuremath{\leftrightarrow}}
\DeclareUnicodeCharacter{21D2}{\ensuremath{\Rightarrow}}
\DeclareUnicodeCharacter{21A6}{\ensuremath{\mapsto}}
\DeclareUnicodeCharacter{2212}{\ensuremath{-}}
\DeclareUnicodeCharacter{2264}{\ensuremath{\leq}}
\DeclareUnicodeCharacter{2265}{\ensuremath{\geq}}
\DeclareUnicodeCharacter{2260}{\ensuremath{\neq}}
\DeclareUnicodeCharacter{2113}{\ensuremath{\ell}}
\DeclareUnicodeCharacter{03A3}{\ensuremath{\Sigma}}
\DeclareUnicodeCharacter{03A6}{\ensuremath{\Phi}}
\DeclareUnicodeCharacter{0393}{\ensuremath{\Gamma}}
\DeclareUnicodeCharacter{221E}{\ensuremath{\infty}}
\DeclareUnicodeCharacter{2208}{\ensuremath{\in}}
\DeclareUnicodeCharacter{2209}{\ensuremath{\notin}}
\DeclareUnicodeCharacter{211A}{\ensuremath{\mathbb{Q}}}
\DeclareUnicodeCharacter{211D}{\ensuremath{\mathbb{R}}}
\DeclareUnicodeCharacter{2124}{\ensuremath{\mathbb{Z}}}
\DeclareUnicodeCharacter{2115}{\ensuremath{\mathbb{N}}}
\DeclareUnicodeCharacter{2200}{\ensuremath{\forall}}
\DeclareUnicodeCharacter{2203}{\ensuremath{\exists}}
\DeclareUnicodeCharacter{2286}{\ensuremath{\subseteq}}
\DeclareUnicodeCharacter{00D7}{\ensuremath{\times}}
\DeclareUnicodeCharacter{03B4}{\ensuremath{\delta}}
\DeclareUnicodeCharacter{2202}{\ensuremath{\partial}}
\DeclareUnicodeCharacter{03B2}{\ensuremath{\beta}}
\DeclareUnicodeCharacter{03B1}{\ensuremath{\alpha}}
\DeclareUnicodeCharacter{03C6}{\ensuremath{\phi}}
\DeclareUnicodeCharacter{03B3}{\ensuremath{\gamma}}
\DeclareUnicodeCharacter{00B9}{\ensuremath{^1}}
\DeclareUnicodeCharacter{00B2}{\ensuremath{^2}}
\DeclareUnicodeCharacter{2070}{\ensuremath{^0}}
\DeclareUnicodeCharacter{2080}{\ensuremath{_0}}
\DeclareUnicodeCharacter{2081}{\ensuremath{_1}}
\DeclareUnicodeCharacter{2082}{\ensuremath{_2}}
\fi
\usepackage{bookmark}
\IfFileExists{xurl.sty}{\usepackage{xurl}}{} % add URL line breaks if available
\urlstyle{same}
\hypersetup{
  pdftitle={The Symbol Game: A First Game with S(n), MSS-N, and Basic Packets},
  pdfauthor={Jeremy H. Carroll},
  colorlinks=true,
  linkcolor={Maroon},
  filecolor={Maroon},
  citecolor={Blue},
  urlcolor={Blue},
  pdfcreator={LaTeX via pandoc}}

\title{The Symbol Game: A First Game with S(n), MSS-N, and Basic
Packets}
\author{Jeremy H. Carroll}
\date{2026-10-02 · Version 0.1.0 · Introductory preprint · Author-review
edition}

\begin{document}
\makeatletter
\begin{center}
{\LARGE \@title\par}
\vspace{0.8em}
{\normalsize \@author\par}
\vspace{0.5em}
{\small \@date\par}
\end{center}
\makeatother

\textbf{Code baseline:}
\nolinkurl{3f3399766108adc4c4e7615dce2573a57be57dd6}

\subsection{Abstract}\label{abstract}

The Symbol Game separates proposing a change from checking that change.
This introduction uses only the Genesis stages S(n), the Minimal Symbol
System atlas MSS-N, and basic packets. A reader can play a first game by
declaring a small collection of marks, requesting a distinction,
proposing a relation, and asking an independent checker to test a stated
rule. A packet records the proposal; a receipt records the check and its
limits. We distinguish a Genesis stage from an atlas entry and from an
execution of the game. Eight gradual examples teach ground counts,
observational separation, compact notation, consistent renaming, typed
prompts and answers, scoped returned packets, bounded sessions, and
extension round trips. Self-contained code and a Jupyter workbook make
each example replayable without claiming to implement Genesis or the
full game. The appendix previews the advanced mathematical and practical
routes, which require additional definitions, evidence, and checkers.

\subsection{1. Start with a question}\label{start-with-a-question}

Imagine two marks, a and b, on a card. You retain both marks but have
not yet said how a reader can distinguish them. Your task is to supply a
distinction and make its checking rule visible. A player proposes a
relation. A referee checks the relation against the marks and the agreed
rule. The referee writes what passed, what failed, and what was not
tested.

That is enough to begin the Symbol Game. A player can be a person, a
program, or a language model. A referee can be another person following
a finite rule or a deterministic program. A program is useful when the
rule is precise enough to replay. Eloquence does not replace a check,
and an accepted packet does not by itself establish every claim in the
player's explanation.

There are indefinitely many possible instances, representations, goals,
and strategies. The first lesson needs only one finite instance. The
advanced game becomes demanding once its checking rules involve serious
mathematics, software, experiments, or combinations of them.

\subsection{2. Three objects to keep
distinct}\label{three-objects-to-keep-distinct}

\textbf{S(n)} names a Genesis derivation stage. New structure must be
justified by an unresolved Genesis obstruction. An external mathematical
interpretation may describe a stage; it cannot create the next stage.
This paper summarizes the repository's core stages and does not supply a
new derivation.

\textbf{MSS-N} names an entry in an externally designed atlas of
representations and checking obligations. Its index labels a
specification. A listed next entry is not a proved transition to that
entry. A proposed transition and a checked, licensed transition are
different records.

\textbf{The Symbol Game} is the proposal-and-check workflow applied to a
declared instance. It can examine a Genesis claim or an MSS obligation,
provided the relevant representation and checking rule are supplied.
Running the workflow does not equate their indices. In particular, MSS-2
is not S(2).

Here is the small part of the Genesis core needed to orient a beginner:

\begin{longtable}[]{@{}
  >{\raggedright\arraybackslash}p{(\linewidth - 2\tabcolsep) * \real{0.5000}}
  >{\raggedright\arraybackslash}p{(\linewidth - 2\tabcolsep) * \real{0.5000}}@{}}
\toprule\noalign{}
\begin{minipage}[b]{\linewidth}\raggedright
Stage
\end{minipage} & \begin{minipage}[b]{\linewidth}\raggedright
Reading for this introduction
\end{minipage} \\
\midrule\noalign{}
\endhead
\bottomrule\noalign{}
\endlastfoot
S(0) & Ground: the presentation \passthrough{\lstinline!<>!} names an
empty carrier. It contains no admitted mark. \\
S(1) & An externally supplied mark may be admitted. Bare supply alone
changes no derivation. \\
S(2) & Marks are retained, but without distinguishing observations they
remain observationally collapsed. \\
S(3) & A demand for separation requires a witnessed distinction.
Directed cuts belong here, not to S(2). \\
S(4) & The declared finite separation task closes when its required
distinctions are checked. \\
S(5) & External continuation introduces a new demand rather than
automatic progress. \\
S(6) & Extension must preserve an explicit way to return to the
predecessor. \\
\end{longtable}

This table is an orientation, not a replacement for the core definitions
and their proof records. Completion of one finite task does not mean all
possible future tasks have been solved. The number of currently
unresolved obstructions is also different from the persistent history of
checked and refused events.

The corresponding beginner atlas route starts at MSS-0 (ground), MSS-1
(an atomic mark), and MSS-2 (relational pairings). A post-hoc comparison
of an atlas representation with a Genesis state is a crosswalk, not
evidence that either generated the other.

\subsection{3. Agree on the board before
playing}\label{agree-on-the-board-before-playing}

For our first game, the carrier is exactly
\passthrough{\lstinline!a, b!}. The goal is to provide a directed
relation from a to b. The permitted relation must name only those two
marks, use the declared relation label, and never relate a mark to
itself. The direction is part of this example's declared task.
Irreflexivity alone would not prove that every pair is separated.

Agree on three things before proposing a move:

\begin{enumerate}
\def\labelenumi{\arabic{enumi}.}
\tightlist
\item
  Which marks and relation label are admitted?
\item
  What finite property will the referee check?
\item
  How many attempts are available, and what happens when they are
  exhausted?
\end{enumerate}

A failed proposal leaves the accepted board unchanged. A stopped attempt
is recorded as unresolved if no checked answer was obtained. A new mark
requires a new declaration or a separately checked extension; it cannot
be smuggled into a repair of the current board.

\subsection{4. A basic packet and a
receipt}\label{a-basic-packet-and-a-receipt}

A packet is a container for a proposal. We use the following teaching
format:

\begin{lstlisting}
{
  "schema": "symbol-game-basic-packet/v1",
  "carrier": ["a", "b"],
  "relation_label": "separates",
  "goal": ["a", "b"],
  "candidate": ["a", "separates", "b"]
}
\end{lstlisting}

This is a deliberately small teaching schema, not a native runtime wire
contract. The examples \passthrough{\lstinline!<>!},
\passthrough{\lstinline!<a>!}, and \passthrough{\lstinline!<ab>!}
elsewhere in the atlas are presentation forms; the JSON fields here are
a convenient way to write a finite checking exercise. Neither notation
licenses a Genesis extension.

The referee tests the schema, uniqueness of the carrier marks,
membership of each endpoint, the relation label, absence of a
self-relation, and agreement with the declared goal. It returns a
receipt such as:

\begin{lstlisting}
{
  "status": "CHECKED",
  "property": "declared-directed-separation",
  "scope": "symbol-game-basic-packet/v1"
}
\end{lstlisting}

The receipt means that the candidate satisfies this finite teaching
rule. It does not certify the semantics of the English word
``separates,'' establish all Genesis obligations, prove an MSS lift, or
authorize an external action. The packet, the checking rule, and the
receipt should be kept together so that someone else can repeat the
check.

\subsection{5. Play and change one
thing}\label{play-and-change-one-thing}

First predict the referee's result without running it. The packet above
passes. Change its candidate to \passthrough{\lstinline!a separates a!}:
the referee refuses a self-relation. Change the second endpoint to c:
the referee refuses an endpoint outside the declared carrier. Reverse
the candidate to \passthrough{\lstinline!b separates a!}: it is
irreflexive, but does not meet this directed goal. Remove the relation
label: the packet is incomplete.

Each refusal teaches a different rule. Repair the candidate within the
same board and try again. If the desired change concerns the carrier or
goal, declare a different instance and preserve the original receipt. Do
not turn an altered question into an apparent solution of the old one.

The companion command
\passthrough{\lstinline!python code/basic\_packets.py!} prints these
outcomes and exits unsuccessfully if a negative control unexpectedly
passes. Its checker uses only finite strings and lists. It does not call
an engine, provider, or network service and makes no claim about
advanced runtime functionals. The local tests also exercise malformed
inputs and a missing goal. These checks support the teaching example,
not a completeness claim about the full game.

\subsection{6. Reading the compact
marks}\label{reading-the-compact-marks}

The ``hieroglyphics tricks'' are ways of writing a declared structure
compactly. They help a reader see boundaries and reuse a pattern. The
meaning comes from an explicit grammar and task, not from the appearance
of the glyphs.

Start with the opaque whole \passthrough{\lstinline!<>!}. Although its
displayed spelling uses two characters, that does not supply two
admitted objects, a direction, or a relation. In the ground exercise
there is one presentation and zero admitted marks. Next,
\passthrough{\lstinline!<a>!} is a supplied whole with a mark. A
supplied \passthrough{\lstinline!<ab>!} still needs a declared
interpretation before its typography can stand for a directed pair. The
Genesis primitives audit makes this boundary explicit.

For the following atlas-style illustrations, declare that
\passthrough{\lstinline!<ab>!} means the directed pairing from a to b
and that \passthrough{\lstinline!<<ab><>>!} records an obstruction
attached to that pairing. Under this convention:

\begin{longtable}[]{@{}
  >{\raggedright\arraybackslash}p{(\linewidth - 2\tabcolsep) * \real{0.5000}}
  >{\raggedright\arraybackslash}p{(\linewidth - 2\tabcolsep) * \real{0.5000}}@{}}
\toprule\noalign{}
\begin{minipage}[b]{\linewidth}\raggedright
Compact display
\end{minipage} & \begin{minipage}[b]{\linewidth}\raggedright
How to read it in this declared exercise
\end{minipage} \\
\midrule\noalign{}
\endhead
\bottomrule\noalign{}
\endlastfoot
\passthrough{\lstinline!<>!} & Ground presentation or
primitive-difference token, according to the stated register \\
\passthrough{\lstinline!<a>!} & A supplied atomic mark \\
\passthrough{\lstinline!<ab>!} & Directed pairing from a to b, when that
interpretation is supplied \\
\passthrough{\lstinline!<ba>!} & The reversed directed pairing \\
\passthrough{\lstinline!<<ab><ba>>!} & A container showing both declared
directed pairings \\
\passthrough{\lstinline!<aa>!} & A self-pair; refused by the declared
irreflexive relation rule \\
\passthrough{\lstinline!<<ab><>>!} & A declared pairing with an attached
obstruction marker \\
\end{longtable}

These are reading examples from the atlas, not a new universal parser.
The companion uses explicit JSON lists for the actual relation checks.
Its result does not establish that every display above is legal in every
wire dialect. The older MSS curriculum includes a different Grammar-000
register with application and dispatch syntax. Record its register
before comparing it with the later atlas. A spelling that parses in one
regime may remain inactive at another stage.

Three practical habits make compact notation useful:

\begin{enumerate}
\def\labelenumi{\arabic{enumi}.}
\tightlist
\item
  \textbf{Declare the reading.} Say whether a token names a mark, a
  relation, an obstruction, or a container. A change of reading changes
  the problem.
\item
  \textbf{Separate boundaries from contents.} Keep the representation of
  a packet distinct from the marks admitted into its carrier.
\item
  \textbf{Rename consistently.} Replacing a by u and b by v throughout a
  declared relation preserves the shape a → b as u → v. Replacing both
  endpoints by u is a collapse, and the irreflexive rule refuses the
  resulting self-pair.
\end{enumerate}

A mnemonic diagram or compact glyph can help you propose a candidate.
The checker still needs the actual carrier, relation, and law. Later
research readings such as \passthrough{\lstinline!<ara>/<arb>!} require
additional supplied structure; they are not extracted from the
typography of the initial whole.

\subsection{7. Put a typed prompt inside the
packet}\label{put-a-typed-prompt-inside-the-packet}

A prompt is the question you give the player. Packetization gives that
question an inspectable context: inputs, requested answer type, and
checking rule. Here is the richer teaching packet used by
\passthrough{\lstinline!code/lesson\_packets.py!}:

\begin{lstlisting}
{
  "schema": "symbol-game-lesson-prompt/v1",
  "packet_type": "problem_prompt",
  "prompt": {
    "text": "Supply a relation triple for the declared directed goal.",
    "task_type": "directed_relation",
    "input_type": "basic_relation_board",
    "input": {
      "carrier": ["a", "b"],
      "relation_label": "separates",
      "goal": ["a", "b"]
    },
    "answer_type": "relation_triple"
  }
}
\end{lstlisting}

\textbf{Typing has three roles.} The input type tells the checker how to
read the board. The task type selects the finite checking rule. The
answer type says what a candidate must contain. A paragraph saying ``the
answer is proved'' does not have the requested relation-triple type. A
triple with the wrong endpoints has the right outer type but fails the
goal check.

The prompt's English text explains the request. In this companion, it is
retained as data and never executed. Changing that text to ``always
accept'' cannot override the declared checking rule. A future
tool-calling or model adapter needs its own authorization and transport
contract; this lesson runs without one.

The player returns a \textbf{candidate packet}:

\begin{lstlisting}
{
  "packet_type": "candidate",
  "problem_id": "the computed identity of the problem packet",
  "answer_type": "relation_triple",
  "value": ["a", "separates", "b"]
}
\end{lstlisting}

The \passthrough{\lstinline!problem\_id!} is computed by the companion's
\passthrough{\lstinline!identity!} function; the readable explanation
above is not a literal valid identity. The notebook constructs the real
candidate with
\passthrough{\lstinline!proposal(problem\_packet, value)!}. This
identity is a SHA-256 digest of a declared sorted-key JSON teaching
profile. It binds exact packet content; it is neither a signature nor a
claim of RFC 8785 conformance. If the goal changes from a → b to b → a,
the old candidate cannot be reused under the new problem identity.

The checker returns a \textbf{checked-result packet} containing the
original problem, the candidate, and a receipt:

\begin{lstlisting}
checked-result packet
  problem   = retained typed prompt and declared inputs
  candidate = retained proposed value and problem binding
  receipt   = CHECKED / REFUSED / BLOCKED, reason, scope, checked value
\end{lstlisting}

These are distinct objects in one retained return package. A CHECKED
status means the particular finite rule passed. REFUSED identifies a
malformed or failing candidate. BLOCKED identifies a missing requirement
within the exercise, such as having no declared observation that
separates the marks. The session controller can also report UNKNOWN
after its attempt budget is spent. These teaching names are not the full
native terminal enumeration.

\subsection{8. A ladder of worked
examples}\label{a-ladder-of-worked-examples}

The examples below adapt useful ideas from the three curriculum strands.
Each has a prediction, a small calculation or check, and a changed-input
control. They are selected finite teaching models. They do not reproduce
every source lesson, implement the complete Genesis derivation engine,
or inherit all historical claims made in older curriculum prose.

\subsubsection{Example 1: one presentation, zero admitted
marks}\label{example-1-one-presentation-zero-admitted-marks}

From the S(0) ground lesson, use:

\begin{lstlisting}[language=Python]
from lesson_packets import examples, proposal, evaluate

ground, collapsed, distinguished, relation = examples()
answer = {"presentation_count": 1, "carrier_count": 0}
returned = evaluate(ground, proposal(ground, answer))
print(returned["receipt"])
\end{lstlisting}

Predict CHECKED. Then change \passthrough{\lstinline!carrier\_count!} to
1. The result is \passthrough{\lstinline!COUNT\_MISMATCH!}. This catches
a common mistake: counting the visible presentation as an admitted
carrier element. The question concerns this explicit ground board; it
does not derive a physical vacuum.

\subsubsection{Example 2: two names, no separating
observation}\label{example-2-two-names-no-separating-observation}

From S(2) retention and collapse, retain marks a and b and supply a
single observation row \passthrough{\lstinline!same, same!}. Both marks
have the same declared observation. Distinct names alone do not witness
observational separation.

\begin{lstlisting}[language=Python]
returned = evaluate(collapsed, proposal(collapsed, 0))
print(returned["receipt"]["reason"])
# NO_DECLARED_SEPARATOR
\end{lstlisting}

The result is BLOCKED. It concerns the finite supplied observation
table, not all possible observations or a proof that separation is
impossible.

\subsubsection{Example 3: supply and check a
separator}\label{example-3-supply-and-check-a-separator}

The next exercise supplies a second observation row
\passthrough{\lstinline!left, right!}. The candidate answer is its index
1:

\begin{lstlisting}[language=Python]
returned = evaluate(distinguished, proposal(distinguished, 1))
print(returned["receipt"]["checked_value"])
# 1
\end{lstlisting}

Index 0 still fails because its row is
\passthrough{\lstinline!same, same!}. This teaches the S(3) separation
demand using a declared observation model. Supplying the new row is an
explicit change of input; the previous BLOCKED receipt remains a record
of the previous board. The teaching model itself licenses no Genesis
advance.

\subsubsection{Example 4: a relation witness and its
law}\label{example-4-a-relation-witness-and-its-law}

From the MSS relational examples and Genesis primitive-difference
lessons, require a supplied directed witness and an irreflexive rule:

\begin{lstlisting}[language=Python]
valid = evaluate(relation, proposal(relation, ["a", "separates", "b"]))
invalid = evaluate(relation, proposal(relation, ["a", "separates", "a"]))
print(valid["receipt"]["status"], invalid["receipt"]["reason"])
# CHECKED SELF_RELATION
\end{lstlisting}

The relation, witness, and law are supplied inputs. The endpoint
separation comes from checking those inputs. Writing a more elaborate
glyph does not supply the missing witness. Reversing the triple gives
another valid outer shape but fails the particular directed goal.

\subsubsection{Example 5: change names while preserving
structure}\label{example-5-change-names-while-preserving-structure}

From the Genesis renaming lesson, rename the entire carrier
consistently:

\begin{lstlisting}[language=Python]
from lesson_packets import rename_edges

marks, edges = rename_edges(["a", "b"], [("a", "b")],
                            {"a": "u", "b": "v"})
print(marks, edges)
# ['u', 'v'] [('u', 'v')]
\end{lstlisting}

The helper requires a bijection on the declared carrier. Change the map
to \passthrough{\lstinline!a → u, b → u!} and it raises a refusal error.
This establishes preservation of this finite directed edge under this
renaming; it does not implement a canonicalizer for every native Genesis
regime.

\subsubsection{Example 6: return a useful refusal, repair, and
replay}\label{example-6-return-a-useful-refusal-repair-and-replay}

Ask the relation checker about \passthrough{\lstinline!a separates a!}.
The return packet preserves both the request and the refused answer,
with reason \passthrough{\lstinline!SELF\_RELATION!} and no checked
value. Repair the answer to \passthrough{\lstinline!a separates b!}
under the same prompt. The returned checked value is the relation
triple, which a caller can consume only after checking the receipt and
its scope.

\begin{lstlisting}[language=Python]
from lesson_packets import replay_return

assert replay_return(valid)
changed = evaluate(relation, proposal(relation, ["a", "separates", "a"]))
changed["receipt"]["status"] = "CHECKED"
assert not replay_return(changed)
\end{lstlisting}

Replay checks the retained problem and candidate again. Relabeling a
refusal cannot turn it into a passing check. This is a local consistency
test, not producer authentication: an attacker who replaces every object
can construct a different self-consistent exercise. Preserve the
expected problem identity when using a return from an untrusted source.

\subsubsection{Example 7: a finite attempt
budget}\label{example-7-a-finite-attempt-budget}

A strategy proposes a self-pair and then a valid directed pair. With one
attempt available, it stops UNKNOWN. With two attempts, the second
candidate passes, and the first refusal is still retained:

\begin{lstlisting}[language=Python]
from lesson_packets import bounded_session

answers = [["a", "separates", "a"], ["a", "separates", "b"]]
print(bounded_session(relation, answers, fuel=1)["status"])
print(bounded_session(relation, answers, fuel=2)["status"])
# UNKNOWN
# CHECKED
\end{lstlisting}

This explains the value of bounded sessions and honest stopping
conditions. A longer budget helped this particular scripted sequence;
the example makes no general claim about convergence, cost savings, or
strategy quality.

\subsubsection{Example 8: keep a way back when
extending}\label{example-8-keep-a-way-back-when-extending}

From the S(6) descent and MSS up/down lessons, retain an explicit
predecessor representation. In a small declared illustration, encode old
marks as \passthrough{\lstinline!("old", mark)!}. The down-map recovers
an admitted old mark; a new term returns
\passthrough{\lstinline!NON\_DESCENT!}:

\begin{lstlisting}[language=Python]
old_marks = {"a", "b"}
def up(mark):
    return ("old", mark)
def down(term):
    return term[1] if term[0] == "old" and term[1] in old_marks else "NON_DESCENT"

assert all(down(up(mark)) == mark for mark in old_marks)
assert down(("new", "c")) == "NON_DESCENT"
\end{lstlisting}

This demonstrates a finite round-trip obligation. It supplies neither a
new Genesis stage nor a licensed MSS lift. A real extension needs its
own source regime, maps, preservation properties, and licensing record.

\subsection{9. The possible stages of a play
session}\label{the-possible-stages-of-a-play-session}

A play session has a lifecycle. Its workflow stages are separate from
the S(n) derivation stages and MSS-N atlas indices.

\begin{longtable}[]{@{}
  >{\raggedright\arraybackslash}p{(\linewidth - 4\tabcolsep) * \real{0.3333}}
  >{\raggedright\arraybackslash}p{(\linewidth - 4\tabcolsep) * \real{0.3333}}
  >{\raggedright\arraybackslash}p{(\linewidth - 4\tabcolsep) * \real{0.3333}}@{}}
\toprule\noalign{}
\begin{minipage}[b]{\linewidth}\raggedright
Session stage
\end{minipage} & \begin{minipage}[b]{\linewidth}\raggedright
What you do
\end{minipage} & \begin{minipage}[b]{\linewidth}\raggedright
What can happen next
\end{minipage} \\
\midrule\noalign{}
\endhead
\bottomrule\noalign{}
\endlastfoot
Declare & Fix the carrier, task, schema, checker, and budget & Begin
when the required inputs and checker exist \\
Packetize & Put the explanation, typed inputs, and requested answer type
in a problem packet & Refuse an invalid prompt; retain a valid
request \\
Propose & Supply a candidate bound to that request & Check it while
treating its explanation as a proposal \\
Check & Validate types and bindings, then the declared finite property &
Return a checked value, a typed refusal, or a blocked obligation \\
Repair & Change a candidate within the same board & Recheck; keep the
earlier refusal \\
Close or stop & Accept the scoped result, or record an unresolved
bounded stop & Replay or compare retained records \\
Request extension & Declare a changed carrier, observation set, or
capability & Check additional obligations; no automatic advance \\
Replay and learn & Recompute the check from the retained objects & Reuse
a scoped result or reproduce a concrete defect \\
\end{longtable}

Some instances close after one candidate. Others require repeated
repairs, branching search, or a new checker. A session may halt because
the answer was checked, because a required capability is absent, or
because its fuel was spent. These are different outcomes. Closing a
local task does not close the whole Symbol Game.

\subsection{10. The value of the returned
packet}\label{the-value-of-the-returned-packet}

A useful return gives the caller something concrete to do. In our
examples, it gives the exact checked count, the index of a separating
observation, or the checked directed relation. A refused return points
to a repairable field. A blocked return identifies an additional input
or obligation. An UNKNOWN session preserves the attempts without making
a conclusion it did not check.

The packet also makes the question portable: another reader can see what
was asked, which value was proposed, and which rule was applied. A
Jupyter notebook can display that sequence; a program can branch on its
typed status. A digest helps detect a changed question. A preserved
refusal makes a changed-input exercise reproducible. None of these alone
proves that a natural-language interpretation is correct.

There are three different values to discuss:

\begin{itemize}
\tightlist
\item
  \textbf{Returned data:} the candidate and any scoped checked value.
\item
  \textbf{Evidence value:} what the reproduced check supports about
  those data.
\item
  \textbf{Practical value:} whether this checked property is useful for
  your real task. That depends on choosing a meaningful carrier,
  question, and checker.
\end{itemize}

For instance, the relation checker gives no useful guarantee about a
medical, legal, or scientific decision merely because its JSON is valid.
A domain task needs the right domain predicate and evidence. This is why
later papers introduce stronger certificates and careful limits.

\subsection{11. What you have learned}\label{what-you-have-learned}

You can now declare a finite board, name its goal, propose a packet,
predict the referee's check, and interpret a scoped receipt. You can
also distinguish an S(n) derivation from an MSS-N specification. These
are sufficient entry skills. Advanced algebra or optimization is
unnecessary for this first game.

The next step is to choose a useful play route. Begin with changed-input
checks and packet repair. Move to graph repair, proof checking, or
experimental comparison only after reading the additional rules for that
route. A finite check can be strong evidence for its declared property
while leaving broader questions open.

\subsection{Mechanism inventory}\label{mechanism-inventory}

The manuscript is introductory exposition. The companion includes finite
schema-and-relation checks, typed prompt and answer bindings,
observation-table checks, bijective renaming, bounded sessions, and
returned-packet replay with positive and negative controls. The S(n)
table summarizes existing core documentation; the MSS descriptions
summarize atlas specifications. Neither companion code nor exposition
derives new Genesis stages, implements the native host, or proves an
advanced lift.

\subsection{A practical reading and notebook
path}\label{a-practical-reading-and-notebook-path}

Start with the
\href{https://github.com/dbccpage/the_symbol_game/blob/3f3399766108adc4c4e7615dce2573a57be57dd6/omega_engine_atlas/coda_atlas/omega_engine_symbol_game/000_pfp_introduction/jupyter_notebooks/000_from_marks_to_returned_packets.ipynb}{Paper
000 workbook}. It executes the eight examples above, prints actual typed
prompt, candidate, and return packets, and asks you to predict
changed-input outcomes. Run it from the package root or its notebook
directory. All inputs are local and all checks have a finite teaching
scope.

The Landau companion provides a learner guide, an examination, and a
separate worked key in
\nolinkurl{_lev_landau_theoretical_minimum_quiz/}. Its core questions
follow the eight examples. Four optional final challenges supply
additional finite definitions and do not license a Genesis or MSS
transition.

Continue with Paper 001 for a stronger worked certificate and Paper 005
for practical session management. Papers 002 and 008 develop
representations, search, and extension obligations. The early papers
should be read as a progressive route: understand the question and
scoped return first, then add the mathematics needed for a stronger
claim.

\begin{longtable}[]{@{}
  >{\raggedright\arraybackslash}p{(\linewidth - 4\tabcolsep) * \real{0.3333}}
  >{\raggedright\arraybackslash}p{(\linewidth - 4\tabcolsep) * \real{0.3333}}
  >{\raggedright\arraybackslash}p{(\linewidth - 4\tabcolsep) * \real{0.3333}}@{}}
\toprule\noalign{}
\begin{minipage}[b]{\linewidth}\raggedright
Curriculum source
\end{minipage} & \begin{minipage}[b]{\linewidth}\raggedright
Idea used here
\end{minipage} & \begin{minipage}[b]{\linewidth}\raggedright
Scope retained
\end{minipage} \\
\midrule\noalign{}
\endhead
\bottomrule\noalign{}
\endlastfoot
S(n), \passthrough{\lstinline!000\_sn\_core!}, lessons 1, 3, 4, and 7 &
Ground counts, collapsed profiles, separation, descent & Finite teaching
representations of core concepts; no new derivation \\
MSS-N, \passthrough{\lstinline!600\_mssn\_frontier!}, early ground,
pairing, and up/down lessons & Syntax versus stage activation, declared
pairings, round-trip checks & Registers must be identified; an example
does not license ascent \\
Genesis, \passthrough{\lstinline!985\_genesis\_primitives!}, lessons 2,
4, and 11 & Supplied relations, refusal boundaries, consistent renaming
& Selected audited models, not a complete or faithful extraction of
Genesis \\
Genesis primitives adversarial audit, correction of 2026-10-01 &
Opaque-whole and attribution boundaries & The relation and its law do
not follow from typography \\
\end{longtable}

\subsection{References}\label{references}

The cited curriculum lives in repository
\href{https://github.com/dbccpage/the_symbol_game}{dbccpage/the\_symbol\_game}.
Its conceptual source baseline is commit
\href{https://github.com/dbccpage/the_symbol_game/commit/3f3399766108adc4c4e7615dce2573a57be57dd6}{\nolinkurl{3f3399766108adc4c4e7615dce2573a57be57dd6}}
(the prefix previously written \passthrough{\lstinline!3f3399766!}).
That pin names the documents below. It is separate from any later
candidate checkout used to execute this package, and it is not a digest
of the files shipped here. Exact hashes of the executed package files
are in
\nolinkurl{docs/release_validation_2026_10_02_review_repairs.json}.
These links identify the cited bytes. This deposit does not re-audit
those unbundled documents.

\begin{itemize}
\tightlist
\item
  Carroll, J. H.
  \href{https://github.com/dbccpage/the_symbol_game/tree/3f3399766108adc4c4e7615dce2573a57be57dd6/omega_engine_education/documentation/0_s_n/core}{\emph{S(n)
  core stages S(0)--S(6)}},
  \passthrough{\lstinline!dbccpage/the\_symbol\_game!} at the source
  baseline above. Definitions of supply, retention, separation, and
  extension.
\item
  Carroll, J. H.
  \href{https://github.com/dbccpage/the_symbol_game/tree/3f3399766108adc4c4e7615dce2573a57be57dd6/omega_engine_education/documentation/1_minimal_symbol_system_n_mssn}{\emph{Minimal
  Symbol System N}}, the same repository and baseline. Externally
  designed atlas and separation of indices and licensed lifts.
\item
  Carroll, J. H. \emph{S(n) curriculum},
  \passthrough{\lstinline!000\_sn\_core!}, lessons on ground, retention,
  required separation, and descent. Repository educational material
  under the core stages linked above.
\item
  Carroll, J. H. \emph{MSS-N curriculum},
  \passthrough{\lstinline!600\_mssn\_frontier!}. Early finite
  representations and explicitly distinguished historical registers,
  under the atlas linked above.
\item
  Carroll, J. H.
  \href{https://github.com/dbccpage/the_symbol_game/blob/3f3399766108adc4c4e7615dce2573a57be57dd6/omega_engine_education/documentation/2_genesis/primitives/README.md}{\emph{Genesis
  primitives: audited educational models}}, correction of 2026-10-01,
  and the
  \href{https://github.com/dbccpage/the_symbol_game/tree/3f3399766108adc4c4e7615dce2573a57be57dd6/omega_engine_education/documentation/2_genesis/curriculum/985_genesis_primitives}{\emph{985\_genesis\_primitives}
  curriculum}. The
  \href{https://github.com/dbccpage/the_symbol_game/blob/3f3399766108adc4c4e7615dce2573a57be57dd6/omega_engine_education/documentation/2_genesis/adversarial_verification/README.md}{adversarial
  audit} supersedes blanket extraction and physical-emergence claims in
  earlier lesson prose.
\item
  Carroll, J. H. \emph{Proposal Fidelity Protocol paper collection},
  Papers 001--010, and the separately retained meron benchmark (Paper
  011). Companion studies and bounded evidence for the advanced routes
  below.
\end{itemize}

\subsection{Appendix A. Preview of the full Symbol
Game}\label{appendix-a.-preview-of-the-full-symbol-game}

The full game is a family of declared finite tasks, not one universal
board. Its deeper routes require the mathematics and implementation
appropriate to the property being checked. This appendix previews those
routes without making them prerequisites of the introduction.

\begin{longtable}[]{@{}
  >{\raggedright\arraybackslash}p{(\linewidth - 2\tabcolsep) * \real{0.5000}}
  >{\raggedright\arraybackslash}p{(\linewidth - 2\tabcolsep) * \real{0.5000}}@{}}
\toprule\noalign{}
\begin{minipage}[b]{\linewidth}\raggedright
Route
\end{minipage} & \begin{minipage}[b]{\linewidth}\raggedright
Additional work and its limit
\end{minipage} \\
\midrule\noalign{}
\endhead
\bottomrule\noalign{}
\endlastfoot
Packet iteration & Validate a richer schema, repair a refusal, and
replay the same instance. Schema validity alone does not prove a
mathematical claim. \\
Exact graph repair (001) & Compare rational residuals and, where
supplied, matching lower-bound certificates. Improving a candidate is
different from proving an optimum. \\
Signed presentations (002) & Check signs, gauges, and bounded search
within a fixed presentation. A finite exclusion is not a physical
classification. \\
Failure detectors (003) & Admit a detector against explicit tests and
negative controls. A constructed example is not a measured error
rate. \\
Native checking and custody (004) & Pin bytes, invoke the relevant host,
and preserve outcomes. Hash equality alone is not semantic
correctness. \\
Practical sessions (005) & Manage repairs, refusal, fuel, and
reproducibility. A stopped session may remain unresolved. \\
Rejected-proposal evidence (006) & Retain and replay counterexamples. A
catalog entry does not establish detector efficacy. \\
Teaching (007) & Predict, execute, and vary examples. A successful
exercise does not establish learning transfer. \\
MSS obligations and lifts (008) & Check context maps, transition
records, and retractions. An atlas row does not license its
successor. \\
Adversarial review (009) & Challenge a concrete property with a
counterexample. Consensus does not replace the property's checker. \\
Holonomy and pilots (010) & Check finite transport and interpret
retained experiments. Small or answer-leaking pilots support no general
model ranking. \\
Meron benchmark (011) & Check a finite cancellation fixture and
distinguish disclosed answers from discovery. The retained pilot is not
evidence of unaided discovery. \\
\end{longtable}

Further representations include rewriting systems, cochains, higher
structures, discrete geometry, quantum information, and bounded
computation. Each introduces new obligations. Mathematical realizations
and physical interpretations remain distinct from the Genesis core.
There is no promise here that every route is implemented, that every
atlas entry is attained, or that a common algorithm solves all
instances.

The practical breadth comes from choosing a carrier, a goal, a proposal
strategy, and a checker. The ten learning routes in the educational
manual are a useful editorial selection, not a measured popularity
ranking or a complete enumeration of ways to play.

\end{document}
